% questions.tex —— 1977 年·西藏卷 · 题目内容片段（无答案）
% 由 tools/generate_projects.py 依来源总账自 1977/1977_数学真题_西藏_试题.md 生成，勿手改（改源 MD 后重跑）


\gkpart{数学试题解（西藏卷）}

\begin{gktip}
\textbf{来源说明}：本题面依据固定源 PDF《西藏1977年高考数学试题解》第 1—3 页逐题转录。固定源是题面与解答连续出现的扫描材料（源书页码 146—147，第 3 页为参考题解答续页、未见页码），每道题的题面后均附解答；正式题为第 1—8 题，原卷未标注分值。第五题之后原卷注明“以下三题任选做二题”，即第六、七、八题任选做二题，本卷按原顺序全部转录。末尾“参考题”两题，源卷注明“不作为录取普通班学生的标准，不评分”，不计入题量。配套答案与解析：答案源自同一解答材料并经独立复算，\textbf{非官方标准答案}。
\end{gktip}

\par\medskip

\begin{question}
\gkti{1}计算：

\nopagebreak[4]
\[
(4.8-3.7)^2-\sqrt[3]{-27}\times\left(5-1\div5\right)\times\left(-\frac58\right)
\]

\end{question}
\begin{question}
\gkti{2}求分式方程

\nopagebreak[4]
\[
\frac{1-x}{x}=x
\]

在 $0$ 到 $1$ 范围内精确到小数三位的根。

\end{question}
\begin{question}
\gkti{3}已知 $\operatorname{Log}_{10}2=0.3010$，$\operatorname{Loe}_{10}5=0.6990$，试计算

\nopagebreak[4]
\[
(2\operatorname{Log}_{10}4)(\operatorname{Log}_25)+\operatorname{Log}_{10}\frac52
\]

的值（精确到小数三位）。

\begin{gktip}
\textbf{源注}：原卷第二个已知条件印作“$\operatorname{Loe}_{10}5$”，系“$\operatorname{Log}_{10}5$”的排印之误（其数值 $0.6990$ 与源解答所用均为 $\operatorname{Log}_{10}5$）。题面照录，计算按 $\operatorname{Log}_{10}5$。
\end{gktip}

\end{question}
\begin{question}
\gkti{4}利用特殊角三角比，求 $\sin75^\circ$ 的值。

\end{question}
\begin{question}
\gkti{5}证明恒等式：

\nopagebreak[4]
\[
(\cos\alpha-\cos\beta)^2+(\sin\alpha-\sin\beta)^2=2-2\cos(\alpha-\beta).
\]

\end{question}
\begin{question}
\gkti{6}红旗工厂要围成一个周长为 $120$ 米的矩形场地可以存放原料。问矩形的长和宽应各取多少米，才能使存放场地的面积最大？并说明这最大面积的场地是什么形状？

\begin{gktip}
\textbf{选题说明}：源卷在第五题之后、第六题之前注明“以下三题任选做二题”，指第六、七、八三题任选做二题；本卷按原顺序全部转录。
\end{gktip}

\end{question}
\begin{question}
\gkti{7}有一正四棱台，上底边长是 $40$ 厘米，下底边长是 $52$ 厘米，高是 $8$ 厘米。求：

（1） 它的下底面的面积是多少？

（2） 它的侧面积是多少？

（3） 它的容积是多少？

\end{question}
\begin{question}
\gkti{8}已知直线 $L_2$ 经过点 $(7,1)$ 并且与直线 $L_1$：$y=\dfrac13x+2$ 垂直，求直线 $L_2$ 的方程。

\Needspace{5\baselineskip}
\textbf{参考题 1．} 求证：

\nopagebreak[4]
\[
1^3+2^3+3^3+\cdots+n^3=\left[\frac{n(n+1)}{2}\right]^2.
\]

\Needspace{5\baselineskip}
\textbf{参考题 2．} 计算定积分的值：

\nopagebreak[4]
\[
\int_0^{\frac{\pi}{4}}\cos2x\,dx.
\]

\end{question}